\chapter{環境認識と地形情報取得}

\section{自己位置推定の概要}

自己位置推定は、ロボットが現在どこにいるのかを推定する技術である。人が部屋の中を移動するときは、視覚情報などを頼りに自身が現在部屋のどこにいるか認識することで、目的地までのルートを決め移動している。当然、目をつむって部屋を移動するのは難しい。このように、移動するためにはその環境において自身の位置を認識することが必須となる。ロボットも同じであり、カメラや測距センサーを用いることで周囲の環境情報を認識し、ロボットの位置を推定している。

自己位置推定の精度は自律移動の精度に大きく直結する。目的地が分かっていても、ロボットの推定自己位置がずれていれば、システム上で目的地に到達しても実際は大きくずれていることがあり得る。特に、不整地環境では、車輪の滑りや地形の起伏により、オドメトリの誤差が蓄積しやすい。そのため、複数のセンサーを組み合わせたセンサーフュージョンにより、高精度な自己位置推定を実現することが重要である。

通常、ロボットの自己位置推定は、事前に周囲の環境地図が分かっている状態で、その地図を基に自己位置推定を行う場合と、事前の環境地図を持たずに未知の環境において地図を作りながら自己位置推定を行う場合がある。本研究では、不整地環境という非常に広い環境で自律移動を行うため、事前に地図を作成するには、環境をくまなく探索する必要がある。また、天候や季節による環境変化も考慮しなければならない。よって、事前に得られた環境地図を基に自己位置推定を行うことは難しく、労力もかかる。そこで、自己位置推定と環境地図を同時に行うSLAM（Simultaneous Localization and Mapping）技術を用いてリアルタイムに地図を作りながら自律移動を行う。

\section{自己位置推定システム}

本研究では、RTK-GPS、wheelオドメトリ、IMUの各オドメトリ情報を統合したセンサーフュージョンにより、高精度な自己位置推定を実現した。先行研究では、これらのセンサーを比較して最適な自己位置推定手法の検討を行い、その結果、GPSオドメトリを用いた自己位置推定がもっとも高精度であることを確認した。そこで、RTK-GPSを主として自己位置推定を行い、wheelオドメトリとIMUを補助的に用いることでGPSオドメトリを補正し、より高精度な自己位置推定を実現した。

自己位置推定は、ロボットの位置に加えて姿勢状態も推定する必要がある。通常、GPSではロボットの座標として緯度・経度・高度に関するデータしか得ることができない。一方、wheelオドメトリでは、wheelエンコーダーから得られる情報を使用して、ロボットの移動に関する推定を行う。これにより、ロボットの位置と姿勢についてのオドメトリデータが得られる。しかし、wheelオドメトリだけでは、車輪の滑りや地形の起伏により、誤差が蓄積しやすい。

IMUは、ロボットの加速度と角速度を測定するセンサーである。IMUから得られるデータを積分することで、ロボットの姿勢と速度を推定できる。しかし、IMUも積分誤差により、長時間の使用では誤差が蓄積する。そこで、これらのセンサーを統合したセンサーフュージョンにより、各センサーの長所を活かし、短所を補う自己位置推定を実現した。

センサーフュージョンには、拡張カルマンフィルタ（Extended Kalman Filter, EKF）を用いる。EKFは、非線形システムにおける状態推定手法であり、複数のセンサーデータを統合して、最適な状態推定を行うことができる。EKFは、予測ステップと更新ステップから構成される。予測ステップでは、システムモデルに基づいて状態を予測し、更新ステップでは、センサーデータを用いて予測を補正する。これにより、高精度な自己位置推定が実現される。

\section{RTAB-Mapによる環境地図生成}

RTAB-Map（Real-Time Appearance-Based Mapping）は、RGB-DカメラやLiDARから得られる視覚情報と3次元点群データを用いて、リアルタイムで環境地図を生成するSLAMライブラリである。RTAB-Mapは、視覚的な特徴点マッチングとループクロージャ検出を組み合わせることで、大規模環境においても長時間にわたり安定して動作する。

RTAB-Mapの主な特徴は、メモリ管理機構を持つことである。通常のSLAMでは、環境地図が大きくなるにつれて、メモリ使用量が増加し、計算コストも増大する。RTAB-Mapは、Working Memory（WM）とLong-Term Memory（LTM）の2階層のメモリ構造を持ち、古いデータを選択的にLTMに移動することで、WMのサイズを一定に保つ。これにより、長時間の運用でもメモリ使用量を一定に保ち、リアルタイム性を維持できる。

ループクロージャ検出は、ロボットが以前訪れた場所に戻ったことを検出する技術である。ループクロージャが検出されると、累積誤差を修正し、地図の整合性を向上させることができる。RTAB-Mapは、Bag-of-Words（BoW）手法を用いて、高速にループクロージャ候補を検索する。BoWは、画像の特徴量を単語に変換し、文書検索の技術を応用して類似画像を検索する手法である。これにより、大規模環境でも高速にループクロージャを検出できる。

本研究では、RTAB-MapにRGB-DカメラとLiDARの両方のセンサーデータを入力し、3次元点群地図を生成する。RGB-Dカメラは、色情報と深度情報を同時に取得できるため、視覚的な特徴点マッチングに有効である。LiDARは、広範囲の高精度な距離測定が可能であり、屋外環境での環境認識に適している。これらのセンサーを統合することで、より高精度で密な点群地図を生成できる。

\section{点群データからの地形情報抽出}

\subsection{ボクセルグリッドの生成}

RTAB-Mapから得られた点群データは、不規則な3次元点の集合であり、そのままでは経路計画に利用しにくい。そこで、点群データをボクセルグリッドに変換する。ボクセルグリッドは、3次元空間を規則的なグリッドセルに分割したデータ構造であり、各セルに点群データが含まれるかどうかを記録する。

本研究では、水平方向に100×100セル、垂直方向に20層のボクセルグリッドを生成する。各セルの解像度は0.2mであり、実空間では20m×20mの領域に相当する。垂直方向の解像度も0.2mであり、最大4mの高さまでの地形を表現できる。このボクセルグリッドにより、不規則な点群データを規則的なデータ構造に変換し、効率的に地形情報を抽出できる。

ボクセルグリッドの生成手順は以下の通りである。まず、点群データの各点について、その3次元座標から対応するボクセルセルを計算する。次に、各ボクセルセルに含まれる点の数を記録する。最後に、各ボクセルセルの代表点として、含まれる点の平均座標を計算する。この代表点を用いて、地形特性を計算する。

\subsection{地形複雑度の計算}

地形複雑度は、標高差、傾斜角、表面粗さの3つの要素から構成される。これらの要素を統合することで、地形の走行適性を包括的に評価できる。

地形複雑度の重みパラメータ（$\alpha=0.4, \beta=0.4, \gamma=0.2$）は、予備実験により決定した。表面粗さの寄与は傾斜・標高差よりも経路品質への影響が小さいことを確認したため、$\gamma$を低く設定した。これにより、ロボットの走行安定性に最も影響を与える標高差と傾斜角を重視した評価が実現された。

標高差$h_{diff}$は、あるセルとその周囲のセルとの高さの差の最大値である。標高差が大きいほど、そのセルは凹凸が激しく、走行が困難であることを示す。標高差は、各セルの代表点の高さを用いて計算される。

傾斜角$\theta_{slope}$は、地形の傾きを表す。傾斜角は、セル内の点群に平面をフィッティングし、その平面の法線ベクトルと鉛直方向とのなす角として計算される。傾斜角が大きいほど、急斜面であり、走行が困難である。

表面粗さ$r_{rough}$は、地形表面の凹凸の細かさを表す。表面粗さは、セル内の点群の高さのばらつき（標準偏差）として計算される。表面粗さが大きいほど、地形表面が粗く、走行が困難である。

これらの3つの要素を重み付け和により統合し、地形複雑度$T_c$を計算する。地形複雑度は以下の式で定義される。

\begin{equation}
T_c = \alpha \cdot h_{diff} + \beta \cdot \theta_{slope} + \gamma \cdot r_{rough}
\label{eq:3_terrain_complexity}
\end{equation}

ここで、$\alpha$、$\beta$、$\gamma$は重みパラメータであり、各要素の重要度を調整する。これらのパラメータは、予備実験により決定した。表\ref{tab:terrain_weight_tuning}に示すように、$\alpha$と$\beta$を0.3から0.5まで0.1刻みで変化させ、$\gamma$を$1-(\alpha+\beta)$として設定し、hill\_detourシナリオにおける経路品質を評価した。

\begin{table}[htbp]
\centering
\caption{地形複雑度重みパラメータの調整結果（hill\_detourシナリオ）}
\label{tab:terrain_weight_tuning}
\begin{tabular}{ccc|ccc}
\hline
$\alpha$ & $\beta$ & $\gamma$ & 経路長[m] & 点平均地形コスト & 総コスト \\
\hline
0.3 & 0.3 & 0.4 & 24.1 & 2.31 & 55.7 \\
0.4 & 0.4 & 0.2 & 27.1 & 1.95 & 52.8 \\
0.5 & 0.5 & 0.0 & 29.8 & 1.82 & 54.2 \\
\hline
\end{tabular}
\end{table}

評価指標として、経路長と点平均法による地形コストの積（総コスト）を用いた。$\alpha=\beta=0.4$、$\gamma=0.2$の設定が、総コスト52.8で最良であったため、本研究ではこの値を採用した。この設定は、標高差と傾斜角が地形の走行適性に大きく影響する一方、表面粗さの寄与は相対的に小さいという物理的解釈とも整合する。

\subsection{地形コストマップの生成}

地形複雑度を経路計画に利用するため、地形コストマップを生成する。地形コストマップは、各セルの地形複雑度を0から1の範囲に正規化したものである。正規化により、異なる地形要素のスケールの違いを吸収し、統一的に扱うことができる。

地形コスト$c_{terrain}$は、以下の式で計算される。

\begin{equation}
c_{terrain} = \frac{T_c}{T_c^{max}}
\label{eq:3_terrain_cost}
\end{equation}

ここで、$T_c^{max}$は、環境内の全セルにおける地形複雑度の最大値である。地形コストが0に近いほど、平坦で走行しやすい地形であり、1に近いほど、複雑で走行困難な地形であることを示す。

生成された地形コストマップは、経路計画アルゴリズムに入力される。経路計画アルゴリズムは、地形コストマップを参照しながら、経路長と地形品質のバランスを考慮した経路を生成する。

\section{まとめ}

本章では、環境認識と地形情報取得の手法について述べた。自己位置推定には、RTK-GPS、wheelオドメトリ、IMUを統合したセンサーフュージョンを用い、拡張カルマンフィルタにより高精度な位置推定を実現した。環境地図生成には、RTAB-Mapを用いて、リアルタイムで3次元点群地図を生成する。点群データからの地形情報抽出では、ボクセルグリッドに変換し、標高差、傾斜角、表面粗さの3つの要素から地形複雑度を計算する。地形複雑度を正規化して地形コストマップを生成し、経路計画アルゴリズムに提供する。次章では、この地形コストマップを利用したTerrain-Aware A*アルゴリズムについて詳しく述べる。
